%% ============================================================================
\section{Methodology}
\label{ch:methodology}
\addcontentsline{lot}{section}{Methodology}
%% ============================================================================

\subsection{Motivation: when the DataLoader becomes the bottleneck}

Deep learning training on HPC is conventionally treated as a GPU problem.
The profiling results reported in Section~\ref{sec:measured-baseline}, together with
a targeted profiling pass over the baseline pipeline, indicate that for this
workload the GPU is frequently not the constraint.
Four findings motivate \pads{}:

\begin{enumerate}
  \item \textbf{GPU starvation.} With \code{num\_workers=0}, the GPU sits
    idle between batches. It is waiting on the data path, not on the model.
  \item \textbf{Shared filesystem pressure.} Thousands of small image and mask files
    generate high metadata overhead on GPFS. The cost is dominated not by bytes
    transferred but by per-file operations against a shared parallel filesystem.
  \item \textbf{Static configuration.} \code{num\_workers}, \code{prefetch\_factor}
    and \code{pin\_memory} are set once at loader construction and never adapt.
    They cannot respond to a bottleneck that shifts - as it does, when contention
    on a shared filesystem varies over the lifetime of a job.
  \item \textbf{Deterministic access.} Because the sample order is fixed once the
    epoch permutation is drawn, the sequence of future file accesses is known.
    Prefetching is therefore not speculative: it can be exact.
\end{enumerate}

\subsection{\pads{}: system overview}

\pads{} (Profiler-Adaptive Data Staging) is a drop-in replacement for the PyTorch
\code{DataLoader} that observes its own behaviour during a short warm-up window,
classifies the dominant bottleneck from the observed timings, selects a staging
policy accordingly, and then executes that policy with predictive prefetching for
the remainder of the run.

The pipeline has four stages:

\begin{enumerate}
  \item \textbf{Warm-up profiling} - 50--100 batches, instrumented.
  \item \textbf{Bottleneck classification} - I/O, decode, H2D or compute.
  \item \textbf{Policy selection} - loose, shard or stage.
  \item \textbf{Predictive staging} - shard prefetch to node-local scratch.
\end{enumerate}

The design intent is that stages 1--3 cost a bounded, one-off overhead - tens of
seconds on a job measured in hours - and that stage 4 recovers that cost many
times over.
The loader is deliberately conservative by default: if no bottleneck is detected, it
degrades to the loose policy, which is the ordinary PyTorch \code{DataLoader}.

\subsubsection{Stage 1 - Warm-up profiling}

For the first 50--100 batches, \pads{} instruments the data path and records
per-batch timings without altering behaviour.
The signals measured are given in Table~\ref{tab:signals}.

\begin{table}[!ht]
\centering
\begin{tabularx}{\textwidth}{@{}lL@{}}
\toprule
Signal & Meaning \\
\midrule
\Tdata{}   & Time spent waiting for the next batch to become available \\
\Tdecode{} & Image and mask decode time, plus augmentation transforms \\
\Thtod{}   & Host-to-device (CPU to GPU) transfer time \\
\Tgpu{}    & Model forward + backward time \\
GPU utilisation & GPU occupancy over the window \\
Queue fullness  & \code{DataLoader} output-queue depth \\
Shard size      & Whether the next shard can be staged efficiently \\
Scratch availability & Whether node-local scratch has free capacity \\
\bottomrule
\end{tabularx}
\caption{Signals recorded during the warm-up window.}
\label{tab:signals}
\end{table}

The final two are feasibility signals rather than performance signals. They
determine whether a policy is \emph{available}, not whether it is desirable.

\paragraph{Instrumentation.}
The four timers are implemented as follows.
\Tgpu{} and \Thtod{} are measured with \code{torch.cuda.Event} pairs recorded around
the transfer and the forward/backward call respectively, synchronised before
reading, because wall-clock timing of asynchronous CUDA work would otherwise
attribute queueing delay to the wrong stage.
\Tdata{} is measured as wall-clock time around \code{next(iterator)}, which captures
the time the training loop is blocked waiting for a batch - and is therefore the
signal that directly expresses starvation.
\Tdecode{} is instrumented inside \code{\_\_getitem\_\_}, accumulating per-sample
decode and transform time, and is aggregated per batch.

The output queue is sampled periodically to measure how full it is. If the queue remains consistently full, the DataLoader is producing batches faster than the model can consume them, indicating that model computation is the bottleneck. Conversely, if the queue remains consistently empty, the model is consuming batches faster than the DataLoader can produce them, indicating that data loading is the bottleneck
It therefore distinguishes between a genuinely efficient data-loading pipeline and a low \Tdata{} that occurs simply because a slow model gives the DataLoader enough time to prefetch and prepare the next batch.


\paragraph{Aggregation.}
Statistics are taken as the median over the window rather than the mean.
The first batches of a run include CUDA context creation, cuDNN algorithm selection
and worker process start-up, all of which are one-off costs that would inflate a
mean and bias the classifier towards diagnosing an I/O bottleneck that does not
exist in steady state.
The first $w_0 = 10$ batches are discarded before statistics are computed, and the
median is taken over the remainder.

\subsubsection{Stage 2 - Bottleneck classification}

The profiled signals are mapped onto a dominant bottleneck class.
The classification is not a black box: it is a small set of rules over the measured
timings, chosen so that the decision is inspectable and reproducible.

\begin{table}[!ht]
\centering
\begin{tabular}{p{4.6cm}p{4.2cm}p{3.4cm}}
\toprule
Detected condition & Interpretation & Action \\
\midrule
\Tdata{} low, \Tgpu{} high, GPU utilisation high
  & Compute-bound; the pipeline is already keeping up & Loose \\
Many \code{open}/\code{stat}/\code{read} calls; high metadata overhead
  & Many-small-file overhead dominates & Shard \\
Shared filesystem read latency high; scratch available
  & Shared-FS latency dominates and can be avoided & Stage \\
Many-small-file overhead \emph{and} slow shared FS
  & Both bottlenecks present & Shard + Stage \\
\Tdecode{} dominates
  & CPU-side decode is the constraint & Increase workers / cache decoded samples \\
\bottomrule
\end{tabular}
\caption{Bottleneck classification rules.}
\label{tab:classification}
\end{table}

Policies are not mutually exclusive.
Scratch staging and an increased worker count may both be active simultaneously when
filesystem I/O and decode are co-bottlenecks.
The classifier therefore emits a \emph{set} of actions rather than a single label.


\paragraph{Threshold-based policy selection.}
The policy-selection rules use relative timing ratios rather than fixed time thresholds, allowing them to remain applicable across GPUs and systems with different performance characteristics. Writing $\tau = \max(\Tdecode{}, \Tgpu{})$, where $\tau$ represents the larger of data-decoding time and GPU computation time, the implementation used here declares:

\begin{align}
\text{stage} &\quad\text{if}\quad \text{scratch available} \;\wedge\; \Thtod{} > \gamma_s \cdot \max(\tau, \epsilon) \\
\text{shard} &\quad\text{if}\quad \Tdecode{} > \gamma_d \cdot \max(\Thtod{}, \Tgpu{}) \;\vee\; \Tdata{} > \gamma_d \cdot \tau \\
\text{loose} &\quad\text{otherwise}
\end{align}

The \code{stage} policy is selected when node-local scratch storage is available and the host-to-device transfer time, \Thtod{}, exceeds the dominant processing time $\tau$ by a factor of $\gamma_s$. The \code{shard} policy is selected when either the decoding time, \Tdecode{}, significantly exceeds both host-to-device transfer and GPU computation time, or when the observed data-waiting time, \Tdata{}, is large relative to $\tau$. If neither condition is satisfied, the loader retains the default \code{loose} policy.

\paragraph{A note on what the classifier optimises.}
The rules above ask ``is the data path the bottleneck?'', not ``which reader is
fastest?''.
These are different questions, and the distinction matters when interpreting
results.
If \Tdata{} is fully hidden behind \Tgpu{} by worker overlap, a reader that is
intrinsically faster buys no wall-clock improvement, and the classifier correctly
declines to switch.
Section~\ref{sec:ablation-thresholds} examines the sensitivity of the outcome to
$\gamma_s$ and $\gamma_d$.

\subsubsection{Stage 3 - Policy selection}

Three staging policies are defined, summarised in Table~\ref{tab:policies}.

\begin{table}[!ht]
\centering
\begin{tabular}{p{2cm}p{6cm}p{5cm}}
\toprule
Policy & What it does & Best when \\
\midrule
Loose & Reads individual image and mask files directly from GPFS
      & Simple baseline; small or node-local dataset; workload already compute-bound \\
Shard & Reads samples from tar archives, amortising per-file metadata cost over many samples
      & Many-small-file overhead is the problem \\
Stage & Copies future shards to node-local scratch ahead of use, removing the shared filesystem from the critical path
      & Shared filesystem latency is the problem \\
\bottomrule
\end{tabular}
\caption{The three staging policies.}
\label{tab:policies}
\end{table}

Loose is the baseline and the fallback.
Shard attacks metadata overhead by converting thousands of \code{open} calls into a
handful of large sequential reads.
Stage attacks latency and contention by moving data off the shared filesystem
entirely, exploiting the fact that the dataset fits in node-local scratch.
The two compose: shards may be both formed and staged.

\paragraph{Archive construction.}
The shard policy requires archives whose member layout the reader can address by
original filename.
Archives are built once, offline, as uncompressed POSIX tar files with members named
\code{images/<name>.jpg} and \code{masks/<name>.png}.
Compression is deliberately not applied: the payload is already-compressed JPEG and
PNG, so entropy coding yields negligible size reduction while forcing the archive to
be decompressed from the stream start on every random access - precisely the
access pattern a shuffled loader produces.
This is a departure from WebDataset~\cite{Aizman19}, which streams shards
sequentially and therefore does not pay that penalty; the choice here preserves
exact random access at the cost of forgoing compression.

\subsubsection{Stage 4 - Predictive staging}

Because the epoch permutation is fixed at the start of each epoch, the order in which the shards will be consumed is known in advance. \pads{} uses this information to stage future shards before they are required. While shard $i$ is being consumed by the training process, a background staging thread copies a future shard $i+d$ from the shared filesystem to node-local scratch, where $d$ denotes the prefetch depth. The objective is to overlap the transfer from shared storage with model computation so that the required shard is already available locally when the training process reaches it.

The prefetch depth $d$ is not a hyperparameter supplied by the user. Instead, it is derived automatically from runtime measurements. Let $t_{\text{stage}}$ denote the measured time required to copy one shard from the shared filesystem to node-local scratch, and let $t_{\text{consume}}$ denote the measured time required by the training process to consume one shard's worth of samples. The depth is set to
%
\begin{equation}
d = \min!\left(
\left\lceil
\frac{t_{\text{stage}}}{t_{\text{consume}}}
\right\rceil + 1,
; d_{\max}
\right),
\qquad
d_{\max} =
\left\lfloor
\frac{C_{\text{scratch}}}{s_{\text{shard}}}
\right\rfloor .
\end{equation}
%
Here, the ratio $t_{\text{stage}}/t_{\text{consume}}$ estimates the number of shard-consumption intervals required to cover the transfer time of one shard. For example, if staging a shard takes $12$ seconds while consuming a shard takes $5$ seconds, the transfer spans approximately $12/5 = 2.4$ consumption intervals. Rounding this value upward gives three shards, and the additional $+1$ provides one shard of safety margin, resulting in a prefetch depth of four. This additional shard provides protection against variations in $t_{\text{stage}}$, which may occur on a shared filesystem due to changing system load and I/O contention.

The calculated depth is additionally constrained by the available node-local scratch capacity. Let $C_{\text{scratch}}$ denote the free capacity of node-local scratch and $s_{\text{shard}}$ the size of a single shard. The resulting value
%
\[
d_{\max} =
\left\lfloor
\frac{C_{\text{scratch}}}{s_{\text{shard}}}
\right\rfloor
\]
%
represents the maximum number of shards that can be stored in scratch at the same time. The final prefetch depth is therefore the smaller of the depth required to hide the measured staging latency and the number of shards that can physically fit in the available scratch space.

This mechanism allows the prefetch depth to adapt automatically to the observed storage and training performance. If $t_{\text{stage}}$ increases, shards must be requested further in advance and the required prefetch depth increases. Conversely, if $t_{\text{consume}}$ increases, the staging thread has more time to complete a transfer while existing shards are being consumed, reducing the required prefetch depth. The prefetch depth is therefore determined by measured system behaviour rather than by a manually selected value.

This is the mechanism by which RQ4---can prefetch depth be auto-tuned from measured timings?---is answered affirmatively by construction. It also distinguishes \pads{} from WebDataset, where shards are streamed using a fixed, statically configured prefetch depth without determining whether that depth is sufficient to hide the measured transfer latency.

\paragraph{Eviction.}
Once a staged shard has been fully consumed, its local copy is removed from node-local scratch to make space for subsequent shards. Since the shard consumption order is known in advance, \pads{} can determine exactly when a shard is no longer required and evict it immediately after consumption. This maintains at most $d$ resident shards and avoids the need for recency-based cache-management heuristics such as least-recently-used (LRU) eviction.

\subsection{Metrics captured}

Two families of metric are reported.

\paragraph{Pipeline metrics (the object of the optimisation).}
\begin{itemize}
  \item \textbf{GPU utilisation} - mean GPU occupancy over the run, the
    headline indicator of whether starvation has been eliminated.
  \item \textbf{Tile throughput} - samples processed per second; the end-to-end
    figure of merit.
  \item \textbf{Time-to-epoch} - wall-clock seconds per epoch, and total
    time-to-target-IoU.
  \item \textbf{\Tdata{} / \Tdecode{} / \Thtod{} / \Tgpu{} breakdown} - retained
    for the full run, not just warm-up, so that the \emph{shift} in the bottleneck
    under each policy is visible rather than merely its disappearance.
  \item \textbf{Filesystem operation counts} - \code{open}/\code{stat}/\code{read}
    calls issued against GPFS, the direct measure of whether sharding achieved what
    it was intended to achieve (RQ2).
  \item \textbf{Scratch traffic and residency} - bytes staged, staging hit rate,
    peak scratch occupancy.
  \item \textbf{Profiling overhead} - the cost of the warm-up window itself,
    reported explicitly so that the net benefit is honest.
\end{itemize}

\paragraph{Model metrics (the constraint, not the objective).}
\begin{itemize}
  \item IoU on the held-out test set, against the Section~\ref{sec:dataset}
    reference, stating the reduction convention used.
\end{itemize}

The reason for reporting both is that this thesis makes a negative claim about
accuracy as well as a positive claim about throughput.
\pads{} is only useful if it accelerates the pipeline \emph{without} degrading
segmentation quality.
A throughput improvement obtained at the cost of reduced shuffling randomness, and
therefore of a worse model, is not an improvement.
Any observed IoU deviation greater than the $\pm 0.5$ percentage point standard
deviation reported in the reference results must be treated as a regression to be
explained, not as noise.

\subsection{Evaluation plan}
\label{sec:evaluation-plan}

\pads{} is compared against four alternatives, all on identical model, dataset and
hardware, with only the data path varying (Table~\ref{tab:evaluation-plan}).

\begin{table}[!ht]
\centering
\begin{tabularx}{\textwidth}{@{}c L L@{}}
\toprule
\# & Method & What it isolates \\
\midrule
1 & Baseline PyTorch \code{DataLoader} -- loose files, default settings
  & The unmodified starting point \\
\addlinespace[2pt]
2 & Tuned PyTorch \code{DataLoader} -- best \code{num\_workers},
    \code{pin\_memory}, \code{prefetch\_factor}
  & How much is recoverable by manual tuning alone; \pads{} must beat this, not merely the naive baseline \\
\addlinespace[2pt]
3 & WebDataset -- fixed tar-shard streaming, no adaptation
  & Isolates the benefit of sharding from the benefit of adaptation (RQ2 vs RQ3) \\
\addlinespace[2pt]
4 & Scratch baseline -- loose files copied manually to node-local scratch up front
  & Isolates the benefit of staging from the benefit of predicting what to stage; the ceiling for any policy that removes GPFS from the path \\
\addlinespace[2pt]
5 & \pads{} (proposed) -- adaptive staging + predictive prefetch
  & --- \\
\bottomrule
\end{tabularx}
\caption{Evaluation methods and the contrast each provides.}
\label{tab:evaluation-plan}
\end{table}

The comparison set is constructed so that each sub-question is answered by a
specific contrast rather than by the aggregate result:

\begin{itemize}
  \item \textbf{RQ1} is answered by the profiling breakdown across methods 1 and 2
    and across the three architectures - SegFormer, with the lowest \Tgpu{}, is
    the configuration in which the loader bottleneck should appear first.
  \item \textbf{RQ2} is answered by 3 vs 1, and by the filesystem operation counts.
  \item \textbf{RQ3} is answered by 5 vs 3, with 4 as the upper bound.
  \item \textbf{RQ4} is answered by 5 vs 3, and by a sweep of fixed depths in the
    ablation of Section~\ref{ch:ablation} - the auto-tuned depth should land at or
    near the best fixed depth, without having been told it.
\end{itemize}

Method~4 deserves emphasis: it is the honest ceiling.
If manually pre-staging the entire dataset to scratch matches \pads{}, then
adaptivity has bought nothing for this dataset, and the contribution reduces to the
observation that the dataset fits in scratch.
\pads{}'s claim is that it reaches that ceiling without requiring the user to know in
advance that scratch is the answer, and that it continues to work when the dataset
does \emph{not} fit - which is the regime the design is actually aimed at, and
which Section~\ref{ch:ablation} probes by synthetically constraining scratch
capacity.

\subsection{Experimental protocol}

To ensure that pipeline differences are not confounded with training noise, all
comparisons hold the following fixed: random seed, epoch permutation, model
architecture and initialisation, batch size, optimiser and learning rate, and the
augmentation pipeline.
Each configuration is intended to be repeated multiple times and reported as mean
$\pm$ standard deviation, because measurements on a shared filesystem are
non-stationary and a single run is not strong evidence on its own. 
